% Sources:
% [ar5iv.labs.arxiv.org](https://ar5iv.labs.arxiv.org/html/2307.09400)
% [arxiv.org](https://arxiv.org/html/2309.17302)
% [ar5iv.labs.arxiv.org](https://ar5iv.labs.arxiv.org/html/1006.3034)
% [math.stonybrook.edu](https://www.math.stonybrook.edu/%7Eoleg/VersoTeX/Show/Hyperfields/hf.html)

\documentclass[11pt]{article}

\usepackage{amsmath,amssymb,amsthm}
\usepackage{hyperref}
\usepackage{url}

\newtheorem{theorem}{Theorem}
\newtheorem{lemma}{Lemma}

\newcommand{\hplus}{\boxplus}
\newcommand{\cP}{\mathcal P}

\title{A Cubic Zero-Sum Identity in Hyperfields}
\author{}
\date{}

\begin{document}

\maketitle

\section{Conventions and standard operations}

Let \(\mathbb H\) be a hyperfield. We use the standard convention that its
hyperaddition
\[
\hplus\colon \mathbb H\times \mathbb H\longrightarrow \cP^{*}(\mathbb H)
\]
takes values in nonempty subsets and is commutative and associative, has \(0\)
as an identity, has unique additive inverses, and satisfies reversibility.
Moreover, \(\mathbb H^\times=\mathbb H\setminus\{0\}\) is an abelian
multiplicative group, \(0\) is absorbing, and scalar multiplication distributes
over hyperaddition:
\[
a(b\hplus c)=ab\hplus ac.
\]
These are the standard hyperfield axioms and the strong form of scalar
distributivity; see \cite{FactoringHyperfields,TropicalExtensions,Viro}.

For nonempty subsets \(A,B\subseteq\mathbb H\), extend hyperaddition by
\[
A\hplus B:=\bigcup_{\substack{a\in A\\ b\in B}}a\hplus b.
\]
We identify an element \(a\in\mathbb H\) with the singleton \(\{a\}\).
Associativity and commutativity of elementwise hyperaddition imply that finite
hypersums of subsets may be flattened, reassociated, and reordered.

\begin{lemma}[Scalar distributivity over finite hypersums]
\label{lem:scalar-distributivity}
For all \(a,x_1,\ldots,x_n\in\mathbb H\),
\[
a(x_1\hplus\cdots\hplus x_n)
=
ax_1\hplus\cdots\hplus ax_n.
\]
\end{lemma}

\begin{proof}
For \(n=2\), this is the scalar distributive law. The general case follows by
induction and associativity:
\[
a\bigl((x_1\hplus\cdots\hplus x_{n-1})\hplus x_n\bigr)
=
a(x_1\hplus\cdots\hplus x_{n-1})\hplus ax_n
=
ax_1\hplus\cdots\hplus ax_n.
\]
The case \(a=0\) is consistent with absorption and the additive identity, since
both sides equal \(\{0\}\).
\end{proof}

\begin{lemma}[Compatibility of negation and multiplication]
\label{lem:negation-multiplication}
For all \(a,b\in\mathbb H\),
\[
-(ab)=(-a)b=a(-b).
\]
\end{lemma}

\begin{proof}
If \(a=0\) or \(b=0\), all three expressions equal \(0\). Otherwise,
\[
0\in a\hplus(-a).
\]
Scalar multiplication by \(b\), together with
Lemma~\ref{lem:scalar-distributivity}, gives
\[
0\in ba\hplus b(-a).
\]
Thus \(b(-a)\) is an additive inverse of \(ba\). By uniqueness of additive
inverses,
\[
b(-a)=-(ba).
\]
Commutativity of multiplication gives the other equality.
\end{proof}

\begin{lemma}[Finite hypersums of sets containing zero]
\label{lem:zero-in-each}
If \(0\in A_i\) for \(1\leq i\leq n\), then
\[
0\in A_1\hplus\cdots\hplus A_n.
\]
\end{lemma}

\begin{proof}
Since \(0\in A_1\) and \(0\in A_2\),
\[
0\in 0\hplus 0=\{0\}\subseteq A_1\hplus A_2.
\]
Iterating this argument gives the claim.
\end{proof}

\begin{lemma}[Balanced-pair augmentation]
\label{lem:balanced-augmentation}
For every subset \(A\subseteq\mathbb H\) and every \(r\in\mathbb H\),
\[
A\subseteq A\hplus r\hplus(-r).
\]
\end{lemma}

\begin{proof}
The additive-inverse axiom gives
\[
0\in r\hplus(-r).
\]
For every \(a\in A\), the identity axiom and monotonicity of extended
hyperaddition give
\[
\{a\}=a\hplus 0
\subseteq a\hplus\bigl(r\hplus(-r)\bigr)
=a\hplus r\hplus(-r).
\]
Taking the union over \(a\in A\) proves the inclusion.
\end{proof}

\section{The cubic identity}

Define
\[
D:=x\hplus y\hplus z
\]
and
\[
T:=x^3\hplus y^3\hplus z^3
\hplus(-xyz)\hplus(-xyz)\hplus(-xyz).
\]
Also define
\[
S:=x^2D\hplus y^2D\hplus z^2D
\hplus(-xy)D\hplus(-xz)D\hplus(-yz)D.
\]

\begin{lemma}[Expansion and regrouping]
\label{lem:expansion}
With the notation above,
\begin{align*}
S={}&T
\hplus\bigl(x^2y\hplus(-x^2y)\bigr)
\hplus\bigl(xy^2\hplus(-xy^2)\bigr)\\
&\quad
\hplus\bigl(x^2z\hplus(-x^2z)\bigr)
\hplus\bigl(xz^2\hplus(-xz^2)\bigr)\\
&\quad
\hplus\bigl(y^2z\hplus(-y^2z)\bigr)
\hplus\bigl(yz^2\hplus(-yz^2)\bigr).
\end{align*}
In particular, \(T\subseteq S\).
\end{lemma}

\begin{proof}
By Lemma~\ref{lem:scalar-distributivity},
\begin{align*}
x^2D&=x^3\hplus x^2y\hplus x^2z,\\
y^2D&=xy^2\hplus y^3\hplus y^2z,\\
z^2D&=xz^2\hplus yz^2\hplus z^3.
\end{align*}
Using Lemma~\ref{lem:negation-multiplication} as well,
\begin{align*}
(-xy)D&=(-x^2y)\hplus(-xy^2)\hplus(-xyz),\\
(-xz)D&=(-x^2z)\hplus(-xyz)\hplus(-xz^2),\\
(-yz)D&=(-xyz)\hplus(-y^2z)\hplus(-yz^2).
\end{align*}
Flattening these six ternary hypersums and reordering the eighteen terms gives
the displayed equality.

Finally, six applications of Lemma~\ref{lem:balanced-augmentation}, with
\[
r=x^2y,\ xy^2,\ x^2z,\ xz^2,\ y^2z,\ yz^2,
\]
show that \(T\subseteq S\).
\end{proof}

The preceding expansion is the hyperfield form of the ordinary polynomial
identity
\[
(x+y+z)(x^2+y^2+z^2-xy-xz-yz)
=x^3+y^3+z^3-3xyz.
\]
It is important that the argument distributes each scalar \(q\) over \(D\)
separately. It does not assert double distributivity of one hypersum over
another, a property not justified by the basic hyperfield axioms; in general
the corresponding product-of-sums statement is only an inclusion
\cite{Viro,ViroNotes}.

\begin{theorem}
Let \(\mathbb H\) be a hyperfield, and let \(x,y,z\in\mathbb H\). If
\[
0\in x\hplus y\hplus z,
\]
then
\[
0\in x^3\hplus y^3\hplus z^3
\hplus(-xyz)\hplus(-xyz)\hplus(-xyz).
\]
\end{theorem}

\begin{proof}
By hypothesis,
\[
0\in D=x\hplus y\hplus z.
\]
For every
\[
q\in\{x^2,y^2,z^2,-xy,-xz,-yz\},
\]
absorption gives \(q0=0\), and hence
\[
0\in qD.
\]
Therefore, by Lemma~\ref{lem:zero-in-each},
\[
0\in
x^2D\hplus y^2D\hplus z^2D
\hplus(-xy)D\hplus(-xz)D\hplus(-yz)D
=S.
\]
Lemma~\ref{lem:expansion} gives \(T\subseteq S\). Since \(0\in S\), it follows
that \(0\in T\), which is precisely the claimed inclusion.
\end{proof}

\begin{thebibliography}{9}

\bibitem{FactoringHyperfields}
\emph{Factoring multivariate polynomials over hyperfields and the
multivariable Descartes' problem}, arXiv:2307.09400.
Available at
\href{https://ar5iv.labs.arxiv.org/html/2307.09400}{ar5iv.labs.arxiv.org}.

\bibitem{TropicalExtensions}
\emph{Geometry of tropical extensions of hyperfields}, Definition~2.1,
arXiv:2309.17302.
Available at
\href{https://arxiv.org/html/2309.17302}{arxiv.org}.

\bibitem{Viro}
O.~Viro,
\emph{Hyperfields for tropical geometry I. Hyperfields and dequantization},
arXiv:1006.3034, Theorems~4.A and~4.B.
Available at
\href{https://ar5iv.labs.arxiv.org/html/1006.3034}{ar5iv.labs.arxiv.org}.

\bibitem{ViroNotes}
O.~Viro,
\emph{Hyperfields and dequantization}, online exposition, Theorem~4.B.
Available at
\href{https://www.math.stonybrook.edu/\%7Eoleg/VersoTeX/Show/Hyperfields/hf.html}{math.stonybrook.edu}.

\end{thebibliography}

\end{document}